\documentclass[conference]{IEEEtran}
\usepackage{amsmath}

\begin{document}

\title{A Flask-Based Student Performance Prediction Prototype Using Linear Regression}

\author{
\IEEEauthorblockN{[Student Name]}
\IEEEauthorblockA{
[Department], [College]\\
[City, Country]\\
[Email Address]
}
}

\maketitle

\begin{abstract}
This paper describes a small web application that estimates a student's final
mark from study hours, previous marks, and attendance. A linear regression
model is trained on the 15 records in the project's included CSV file and
served through a Flask API. Leave-one-out cross-validation (LOOCV) gives a
mean absolute error (MAE) of 0.640 marks, a root mean squared error (RMSE) of
0.857 marks, and an $R^2$ of 0.998. A mean-prediction baseline gives an MAE of
16.648 marks, an RMSE of 19.563 marks, and an $R^2$ of -0.148. These results
describe only the included sample. Its source is undocumented and it contains
only 15 observations; therefore, the scores do not establish that the model
will generalize to other students. The application is presented as an
educational prototype, not as a validated grading or decision-making system.
\end{abstract}

\begin{IEEEkeywords}
student performance, linear regression, Flask, cross-validation, educational
prototype
\end{IEEEkeywords}

\section{Introduction}
Predicting student performance is a useful way to learn how a dataset, a
machine-learning model, and a web application work together. This project
uses study hours, previous marks, and attendance to estimate final marks.
Its purpose is to demonstrate a basic regression workflow rather than to
assess students or replace an instructor's judgment.

\section{System and Method}
\subsection{Dataset}
The included file, \texttt{students.csv}, contains 15 records and four numeric
columns: study hours, previous marks, attendance, and final marks. The first
three columns are the model features, and final marks is the prediction
target. The data source and collection method are not documented. The results
are therefore treated as exploratory results on demonstration data, not as
evidence from a defined student population.

\subsection{Regression Model}
The model is scikit-learn's ordinary least-squares \texttt{LinearRegression},
which estimates a weighted sum of input features [1]. Fitted on all 15 rows,
the approximate equation is
\begin{equation}
\hat{y} = 9.223 + 2.981h + 0.718m + 0.017a,
\end{equation}
where $h$ is study hours, $m$ is previous marks, and $a$ is attendance.
These fitted coefficients describe this dataset; they do not establish
causation. The application limits the displayed prediction to the range
0--100 and assigns a performance label using fixed thresholds.

\subsection{Web Application}
The browser interface collects the three numeric inputs and sends them to a
Flask \texttt{/predict} endpoint. The server checks that each input is
present, finite, and within its allowed range before making a prediction.
The student's name is displayed in the browser and is not included in the
prediction request.

\section{Evaluation}
LOOCV was used because the dataset is small. Each of the 15 records is held
out once, and the model is trained on the other 14 records for that fold.
Cross-validation makes use of the available records but does not replace
testing on a larger independent dataset [2]. The same procedure was applied
to a baseline that predicts the mean final mark of the training fold. Metrics
were computed over all 15 held-out predictions. The results can be reproduced
by running \texttt{python evaluate\_model.py}.

\begin{table}[htbp]
\caption{LOOCV Results on the Included 15-Record Dataset}
\label{tab:results}
\centering
\begin{tabular}{lrrr}
\hline
Model & MAE & RMSE & $R^2$ \\
\hline
Linear regression & 0.640 & 0.857 & 0.998 \\
Mean baseline & 16.648 & 19.563 & -0.148 \\
\hline
\end{tabular}
\end{table}

\section{Discussion}
On the included records, linear regression has lower prediction errors than
the mean baseline. This comparison is useful for understanding the prototype,
but the very small sample and undocumented data origin limit what can be
concluded. The results may reflect a pattern in these rows and should not be
presented as verified accuracy for students outside this dataset. The three
features also omit many circumstances that may affect learning.

\section{Limitations and Responsible Use}
The dataset contains only 15 records, its source is undocumented, and no
independent test set is available. The model is a simple linear model and the
selected features cannot represent a student's full circumstances. A
prediction should not be used as an official grade or for a high-stakes
decision. A responsible follow-up would document data provenance, obtain
permission for data use, collect a larger and appropriate dataset, protect
student privacy, and evaluate on independent records.

\section{Conclusion}
This mini-project connects a CSV dataset, a linear regression model, a Flask
API, and a browser interface. LOOCV on the included records shows lower errors
than a mean baseline, but the result is limited to this small,
undocumented dataset. The project is a learning prototype; further data
collection and independent evaluation would be necessary before making
real-world claims.

\begin{thebibliography}{00}
\bibitem{sklearn-linear}
Scikit-learn developers, ``Linear Models,'' \emph{Scikit-learn User Guide}.
[Online]. Available: https://scikit-learn.org/stable/modules/linear\_model.html

\bibitem{sklearn-cv}
Scikit-learn developers, ``Cross-validation: evaluating estimator
performance,'' \emph{Scikit-learn User Guide}. [Online]. Available:
https://scikit-learn.org/stable/modules/cross\_validation.html
\end{thebibliography}

\end{document}
